%%%%%%%%%%%%%%%%%%%%%%%%%%%%%%%% Define Article %%%%%%%%%%%%%%%%%%%%%%%%%%%%%%%
\documentclass[14pt]{extarticle}
%%%%%%%%%%%%%%%%%%%%%%%%%%%%%%%%%%%%%%%%%%%%%%%%%%%%%%%%%%%%%%%%%%%%%%%%%%%%%%%

%%%%%%%%%%%%%%%%%%%%%%%%%%%%%%%% Using Packages %%%%%%%%%%%%%%%%%%%%%%%%%%%%%%%
\usepackage[T2A]{fontenc}
\usepackage[utf8]{inputenc}
\usepackage[main=russian, english]{babel}

\usepackage{geometry}
\usepackage{esvect}
\usepackage{bm}
\usepackage{amsmath, amsthm, amssymb}
\usepackage{mathtools}
\usepackage{mathrsfs}
\usepackage{nicematrix}
\usepackage{array}
\usepackage{enumitem}
\usepackage{ytableau}
\usepackage{mathdots}
\usepackage[most]{tcolorbox}

\usepackage{indentfirst}
\usepackage{import}
\usepackage{fancybox, fancyhdr}
\usepackage[colorlinks=true, linkcolor=blue, unicode]{hyperref}
\usepackage{tocloft}
\usepackage{titlesec}

\usepackage{graphics}
\usepackage{pgfplots}
\usepackage{tikz}
%%%%%%%%%%%%%%%%%%%%%%%%%%%%%%%%%%%%%%%%%%%%%%%%%%%%%%%%%%%%%%%%%%%%%%%%%%%%%%%

% Other Settings

%%%%%%%%%%%%%%%%%%%%%%%%%%%%%%%% Page Setting %%%%%%%%%%%%%%%%%%%%%%%%%%%%%%%%%
\geometry{a4paper, top=2cm, right=1.5cm, bottom=2cm, left=1.5cm}
\setlength{\parindent}{20pt}
\setlength{\headheight}{17pt}
%%%%%%%%%%%%%%%%%%%%%%%%%%%%%%%%%%%%%%%%%%%%%%%%%%%%%%%%%%%%%%%%%%%%%%%%%%%%%%%

%%%%%%%%%%%%%%%%%%%%%%%%%% Table of contents Setting %%%%%%%%%%%%%%%%%%%%%%%%%%
\renewcommand{\cfttoctitlefont}{\hfill\Large\bfseries}
\renewcommand{\cftaftertoctitle}{\hfill}

\renewcommand{\cftsecaftersnum}{.}
\renewcommand{\cftsubsecaftersnum}{.}
\renewcommand{\cftsubsubsecaftersnum}{.}

\renewcommand{\cftsecdotsep}{\cftdotsep}
\renewcommand{\cftsubsecdotsep}{\cftdotsep}
\renewcommand{\cftsubsubsecdotsep}{\cftdotsep}
%%%%%%%%%%%%%%%%%%%%%%%%%%%%%%%%%%%%%%%%%%%%%%%%%%%%%%%%%%%%%%%%%%%%%%%%%%%%%%%

%%%%%%%%%%%%%%%%%%%%%%%%%%%%%%%% Title Setting %%%%%%%%%%%%%%%%%%%%%%%%%%%%%%%%
\titleformat{\section}
    {\centering\LARGE\bfseries} % format
    {\thesection. } % label
    {0pt} % sep
    {\LARGE} % size

\titleformat{\subsection}
    {\centering\large\bfseries} % format
    {\thesubsection. } % label
    {0pt} % sep
    {\large} % size

\titleformat{\subsubsection}
    {\centering\normalsize\bfseries} % format
    {\thesubsubsection. } % label
    {0pt} % sep
    {\normalsize} % size
%%%%%%%%%%%%%%%%%%%%%%%%%%%%%%%%%%%%%%%%%%%%%%%%%%%%%%%%%%%%%%%%%%%%%%%%%%%%%%%

%%%%%%%%%%%%%%%%%%%%%%%%%%%%%%% Drawing Settings %%%%%%%%%%%%%%%%%%%%%%%%%%%%%%
\usetikzlibrary{calc, arrows.meta, angles, quotes}
%%%%%%%%%%%%%%%%%%%%%%%%%%%%%%%%%%%%%%%%%%%%%%%%%%%%%%%%%%%%%%%%%%%%%%%%%%%%%%%

%%%%%%%%%%%%%%%%%%%%%%%%%%%%%% Plotting Settings %%%%%%%%%%%%%%%%%%%%%%%%%%%%%%
\usepgfplotslibrary{fillbetween}
\pgfplotsset{compat=newest}
%%%%%%%%%%%%%%%%%%%%%%%%%%%%%%%%%%%%%%%%%%%%%%%%%%%%%%%%%%%%%%%%%%%%%%%%%%%%%%%

%%%%%%%%%%%%%%%%%%%%%%%%%%%%%%%%%%% Colors %%%%%%%%%%%%%%%%%%%%%%%%%%%%%%%%%%%%
\definecolor{mygreen}{HTML}{a5d88a}
\definecolor{myyellow}{HTML}{f5e04c}
\definecolor{myorange}{HTML}{ffb299}
\definecolor{mypink}{HTML}{f0b6e0}
\definecolor{myblue}{HTML}{c8e1f2}
\definecolor{mypurple}{HTML}{b19cd9}
\definecolor{mygray}{HTML}{cccccc}
%%%%%%%%%%%%%%%%%%%%%%%%%%%%%%%%%%%%%%%%%%%%%%%%%%%%%%%%%%%%%%%%%%%%%%%%%%%%%%%

%%%%%%%%%%%%%%%%%%%%%%%%%%%%%%%%%%%%%%%%%%%%%%%%%%%%%%%%%%%%%%%%%%%%%%%%%%%%%%%
%\newtheorem{definition}{Определение}[section]
%\newtheorem{theorem}{Теорема}[section]
%\newtheorem{lemma}{Лемма}[section]
%\newtheorem{proposition}{Предложение}[section]
%\newtheorem*{remark}{Замечание}
%\newtheorem{example}{Пример}[section]
%\newtheorem*{corollary}{Следствие}
\numberwithin{equation}{section}
%\newtheorem{exercise}{Упражнение}[section]
%%%%%%%%%%%%%%%%%%%%%%%%%%%%%%%%%%%%%%%%%%%%%%%%%%%%%%%%%%%%%%%%%%%%%%%%%%%%%%%

%%%%%%%%%%%%%%%%%%%%%%%%%%%%%%%%%% Math Sets %%%%%%%%%%%%%%%%%%%%%%%%%%%%%%%%%%
\newcommand{\BB}{\mathbb{B}}
\newcommand{\NN}{\mathbb{N}}
\newcommand{\ZZ}{\mathbb{Z}}
\newcommand{\QQ}{\mathbb{Q}}
\newcommand{\II}{\mathbb{I}}
\newcommand{\RR}{\mathbb{R}}
\newcommand{\CC}{\mathbb{C}}
%%%%%%%%%%%%%%%%%%%%%%%%%%%%%%%%%%%%%%%%%%%%%%%%%%%%%%%%%%%%%%%%%%%%%%%%%%%%%%%

%%%%%%%%%%%%%%%%%%%%%%%%%%%%%%%%%%%%%%%%%%%%%%%%%%%%%%%%%%%%%%%%%%%%%%%%%%%%%%%
\newcommand{\silentsection}[1]{
    \refstepcounter{section}
    \addcontentsline{toc}{section}{\protect#1}
    \section*{#1}
    \sectionmark{#1}
}

\newcommand{\silentsubsection}[1]{
    \refstepcounter{subsection}
    \addcontentsline{toc}{subsection}{\protect#1}
    \subsection*{#1}
    \subsectionmark{#1}
}

\newcommand{\silentsubsubsection}[1]{
    \refstepcounter{subsubsection}
    \addcontentsline{toc}{subsubsection}{\protect#1}
    \subsubsection*{#1}
    \subsubsectionmark{#1}
}
%%%%%%%%%%%%%%%%%%%%%%%%%%%%%%%%%%%%%%%%%%%%%%%%%%%%%%%%%%%%%%%%%%%%%%%%%%%%%%%

%%%%%%%%%%%%%%%%%%%%%%%%%%%%%%%%%%%%%%%%%%%%%%%%%%%%%%%%%%%%%%%%%%%%%%%%%%%%%%%
\tcbset{
    myboxstyle/.style={
        enhanced,
        breakable,
        fonttitle=\bfseries,
        coltitle=black,
        arc=3mm,
        top=2mm,
        bottom=2mm,
        left=3mm,
        right=3mm,
        boxsep=1mm,
        drop fuzzy shadow={black!30!white}
    }
}
\newtcbtheorem[number within=section]
    {defn}{Определение}
    {myboxstyle,
     colback=mygreen!10,
     colframe=mygreen!60
    }
    {defn}
\newtcbtheorem[use counter from=defn]
    {prop}{Предложение}
    {myboxstyle,
     colback=myblue!10,
     colframe=myblue!90
    }
    {prop}
\newtcbtheorem[use counter from=defn]
    {thrm}{Теорема}
    {myboxstyle,
     colback=myorange!10,
     colframe=myorange!70
    }
    {thrm}
\newtcbtheorem[use counter from=defn]
    {lemma}{Лемма}
    {myboxstyle,
     colback=mypink!10,
     colframe=mypink!60
    }
    {lemma}
\newtcbtheorem[use counter from=defn]
    {coroll}{Следствие}
    {myboxstyle,
     colback=mypink!10,
     colframe=mypink!60
    }
    {coroll}
\newtcbtheorem[use counter from=defn]
    {example}{Пример}
    {myboxstyle,
     colback=myyellow!10,
     colframe=myyellow!60
    }
    {example}
\newtcbtheorem[use counter from=defn]
    {exercise}{Упражнение}
    {myboxstyle,
     colback=mypurple!10,
     colframe=mypurple!60
    }
    {exercise}
\newtcbtheorem[use counter from=defn]
    {remark}{Замечание}
    {myboxstyle,
     colback=mygray!10,
     colframe=mygray!60
    }
    {remark}
%%%%%%%%%%%%%%%%%%%%%%%%%%%%%%%%%%%%%%%%%%%%%%%%%%%%%%%%%%%%%%%%%%%%%%%%%%%%%%%

%%%%%%%%%%%%%%%%%%%%%%%%%%%%%%% Math Operators %%%%%%%%%%%%%%%%%%%%%%%%%%%%%%%%
\DeclareMathOperator{\spec}{spec}
\DeclareMathOperator{\tr}{tr}
%%%%%%%%%%%%%%%%%%%%%%%%%%%%%%%%%%%%%%%%%%%%%%%%%%%%%%%%%%%%%%%%%%%%%%%%%%%%%%%

%%%%%%%%%%%%%%%%%%%%%%%%%%%% Title & Author & Date %%%%%%%%%%%%%%%%%%%%%%%%%%%%
\title{\Huge Задачи к коллоквиуму}
\author{\huge Группа 251}
\date{2025 --- 2026}
%%%%%%%%%%%%%%%%%%%%%%%%%%%%%%%%%%%%%%%%%%%%%%%%%%%%%%%%%%%%%%%%%%%%%%%%%%%%%%%

\begin{document}
    \pagestyle{fancy}
    \maketitle
    \tableofcontents
    \newpage

    \silentsection{Определения}

    \silentsubsection{1}

    \textbf{Условие:}

    Определение того, что множество \( A \) лежит левее множества \( B \),
    разделяющего элемента и формулировка принципа полноты.
    Определение множества вещественных чисел (с помощью \( 11 \) свойств)
    и формулировка принципа полноты для множества вещественных чисел.

    \bigskip \bigskip

    \begin{itemize}
        \item
            \( A \leq B \)

            \( A \) лежит левее множества \( B \), если
            \( \forall a \in A, \ b \in B : a \leq b \)
        \item
            Принцип полноты и разделяющий элемент.

            Множество \( C \) обладает свойством полноты, если
            \begin{gather*}
                \text{для любых непустых} \ A, B \subseteq C : A \leq B
                \\
                \text{существует разделяющий элемент} \ c \in C : \forall a \in A, \ b \in B : a \leq c \leq b
            \end{gather*}
        \item
            Вещественные числа.

            \begin{enumerate}[start=0]
                \item {
                    Замкнутость относительно основных арифметических операций.
                }
                \item {
                    \( \forall a, b, c \in R \ a + (b + c) = (a + b) + c \)
                }
                \item {
                    \( \forall a, b \in R \ a + b = b + a \)
                }
                \item {
                    \( \exists 0 \in R : a + 0 = a \)
                }
                \item {
                    \( \forall a \in R \ \exists b \in R : a + b = 0 \)
                }
                \item {
                    \( \forall a, b, c \in R \ a \cdot (b \cdot c) = (a \cdot b) \cdot c \)
                }
                \item {
                    \( \forall a, b \in R \ a \cdot b = b \cdot a \)
                }
                \item {
                    \( \exists 1 \in R, 1 \neq 0 : 1 \cdot a = a \)
                }
                \item {
                    \( \forall a \in R, a \neq 0 \ \exists b \in R, b \neq 0: a \cdot b = 1 \)
                }
                \item {
                    \( \forall a, b, c \in R \ a \cdot (b + c) = a \cdot b + a \cdot c \)
                }
                \item {
                    \( \forall a, b \in R \ a \leq b \lor b \leq a \)

                    \( a \leq b, \forall c \in R \ a + c \leq b + c \)

                    \( a \leq b, \forall c \in R, c \geq 0 \ a \cdot c \leq b \cdot c \)
                }
                \item {
                    На \( R \) выполнено свойство полноты.
                }
            \end{enumerate}

            Если определить \( \RR \), как бесконечные десятичные дроби (запретив период из девяток),
            то для них выполнен принцип полноты --- существует разделяющая дробь.
    \end{itemize}

    \newpage

    \silentsubsection{2}

    \textbf{Условие:}

    Аксиома Архимеда и следствие из аксиомы Архимеда.
    Система и последовательность вложенных и стягивающихся отрезков и лемма о вложенных отрезках.

    \bigskip \bigskip

    \begin{itemize}
        \item
            Аксиома Архимеда

            Для любого положительного вещественного числа \( a \) существует такое натуральное число \( n \),
            что \( na \geq 1 \).

            Отсюда вытекает полезное следствие:
            \begin{enumerate}
                \item
                    \( \forall x, y \in \RR : x < y \ \exists c \in \QQ : x < c < y \)
                \item
                    \( \forall x, y \in \RR : x < y \ \exists c \in \II : x < c < y \)
            \end{enumerate}
        \item
            Система и последовательность вложенных и стягивающихся отрезков.

            \( [a, b] = \{ x \in \RR \ \vline \ a \leq x \leq b \} \) --- отрезок множества действительных чисел.

            Системой вложенных отрезков называется множество отрезков \( M \),
            в котором для любых \( I_1, \ I_2 \in M \)
            выполнено \( I_1 \subset I_2 \) или \( I_2 \subset I_1 \)

            Если, дополнительно, в \( M \) все отрезки занумерованы,
            и любой отрезок с большим номером содержится в любом отрезке с меньшим номером,
            то множество \( M \) называют последовательностью вложенных отрезков.

            Последовательность вложенных отрезков называется стягивающейся,
            если \( \forall \varepsilon > 0 \) существует отрезок с длиной, меньшей \( \varepsilon \).
        \item
            Лемма о вложенных отрезках.

            Если дана последовательность вложенных отрезков \( M \),
            то \( \exists c : \forall I \in M \ c \in I \)
            (все отрезки \( M \) имеют общий элемент).

            Если при этом \( M \) --- последовательность стягивающих отрезков,
            то общий элемент единственный.
    \end{itemize}

    \newpage

    \silentsubsection{3}

    \textbf{Условие:}

    Определение числовой последовательности и все определения предела последовательности

    \bigskip \bigskip

    \begin{itemize}
        \item
            Числовая последовательность.

            Декартово произведение множеств \( X \) и \( Y \) обозначается \( X \times Y \)
            и является множеством всех упорядоченных пар \( (x, y) \), где \( x \in X, \ y \in Y \).

            (btw упорядоченная пара \( (x, y) \) --- \( \{x, \{ x, y \} \} \))

            Подмножество \( f \) декартова произведения \( X \times Y \) называется функцией \\
            из \( X \) в \( Y \) (\( f : X \rightarrow Y \)), если
            \( \forall x \in X \ \exists ! \ (x, y) \in f \).
            При этом пишут \( f(x) = y \).

            \( f : \NN \to \RR \) называется числовой последовательностью.
            Значение \( f(n) \) обозначается \( a_n \).
            Также последовательность обозначают \( \{ a_n \}_{n = 1}^\infty \)
        \item
            Предел последовательности.

            \( (a, b) = \{ x \in \RR \ \vline \ a < x < b \} \) --- интервал множества действительных чисел.

            Окрестностью точки \( c \) называется любой интервал \( (a, b) \) такой, что \( c \in (a, b) \).

            \( U(c) \) --- окрестность точки \( c \)

            \( U_\varepsilon (c) = (c - \varepsilon, c + \varepsilon) \) --- \( \varepsilon \)-окрестность точки \( c \)

            \begin{enumerate}
                \item
                    Через окрестности:
                    \begin{gather*}
                        \lim\limits_{n \to \infty} a_n = A
                        \Leftrightarrow
                        \forall U(A) \ \exists N : \forall n \geq N \ a_n \in U(A)
                        \\
                        \lim\limits_{n \to \infty} a_n = A
                        \Leftrightarrow
                        \forall \varepsilon > 0 \ \exists N : \forall n \geq N \ a_n \in U_\varepsilon(A)
                    \end{gather*}
                \item
                    Через бесконечно малую (вероятно, ненужно):

                    \( \{ a_n \} - \text{БМП, если} \ \lim\limits_{n \to \infty} a_n = 0 \)

                    \( \lim\limits_{n \to \infty} a_n = A \Leftrightarrow \)
                    существует БМП \( \{ \alpha_n \} \) такая, что \( a_n = A + \alpha_n \).
            \end{enumerate}
    \end{itemize}

    \newpage

    \silentsubsection{4}

    \textbf{Условие:}

    Единственность предела последовательности, определение ограниченной последовательности,
    ограниченность сходящейся последовательности и лемма об отделимости.
    Арифметика пределов, переход к пределу в неравенстве и лемма о зажатом пределе.

    \bigskip \bigskip

    \begin{itemize}
        \item
            Единственность предела последовательности.

            Если последовательность имеет предел, то он единственен.
        \item
            Ограниченная последовательность.

            \( \{ a_n \} \) ограничена, если \( \exists C : \forall n \ |a_n| < C \)
        \item
            Ограниченность сходящейся последовательности.

            \( \{ a_n \} \) сходится \( \Rightarrow \) \( \{ a_n \} \) ограничена.
        \item
            Лемма об отделимости.

            \( \lim\limits_{n \to \infty} a_n = A \) и \( A \neq 0 \)
            \( \Rightarrow \exists N : \forall n > N \ |a_n| > \frac{|A|}{2} \)
        \item
            Арифметика пределов.

            Пусть \( \lim\limits_{n \to \infty} a_n = A \) и \( \lim\limits_{n \to \infty} b_n = B \)

            Тогда:
            \begin{enumerate}
                \item
                    \( \forall \alpha, \beta \in \RR \ \lim\limits_{n \to \infty} (\alpha a_n + \beta b_n) = \alpha A + \beta B \)
                \item
                    \( \lim\limits_{n \to \infty} a_n b_n = AB \)
                \item
                    Если \( \forall n \ b_n \neq 0 \),
                    то \( \lim\limits_{n \to \infty} \frac{a}{b} = \frac{A}{B} \)

                    (на самом деле, благодаря лемме об отделимости, достаточно \( B \neq 0 \))
            \end{enumerate}
        \item
            Предельный переход в неравенстве.

            Пусть \( \lim\limits_{n \to \infty} a_n = A \) и \( \lim\limits_{n \to \infty} b_n = B \),
            причем \( \exists N : \forall n > N \ a_n \leq b_n \).

            Тогда \( A \leq B \).

            \textbf{Важно}, если заменить \( a_n \leq b_n \) на \( a_n < b_n \), то все еще будет \( A \leq B \)
        \item
            Лемма о зажатом пределе.

            Пусть \( \lim\limits_{n \to \infty} a_n = A \) и \( \lim\limits_{n \to \infty} b_n = A \),
            причем \( \exists N : \forall n > N \ a_n \leq c_n \leq b_n \).

            Тогда \( \exists \lim\limits_{n \to \infty} c_n \),
            более того --- \( \lim\limits_{n \to \infty} c_n = A \)
    \end{itemize}

    \newpage

    \silentsubsection{5}

    \textbf{Условие:}

    Определение бесконечно малой последовательности,
    свойство произведения бесконечно малой последовательности на ограниченную
    и определение предела в терминах бесконечно малых последовательностей.
    Определение верхних и нижних граней, а также два определения точной верхних и нижних граней.
    Существование точной верхней и нижней граней ограниченного множеств.

    \bigskip \bigskip

    \begin{itemize}
        \item
            Бесконечно малая последовательность.

            \( \{ a_n \} - \text{БМП, если} \ \lim\limits_{n \to \infty} a_n = 0 \)

            Если \( \{ a_n \} \) --- бесконечно малая, а \( \{ b_n \} \) --- ограниченная,
            то \( \{ a_n b_n \} \) --- бесконечно малая.

            \( \lim\limits_{n \to \infty} a_n = A \Leftrightarrow \)
            существует БМП \( \{ \alpha_n \} \) такая, что \( a_n = A + \alpha_n \).
        \item
            Грани.

            Пусть \( A \subset \RR, A \neq \varnothing \).

            \begin{enumerate}
                \item
                    Если существует \( \exists C \in \RR \ \forall a \in A \ a \leq C \),
                    то \( C \) называется верхней гранью множества \( A \).

                    Наименьшая из верхних граней множества \( A \) называется точной верхней гранью
                    и обозначается \( \sup \) (супремум).

                    \( C = \sup A \), если:
                    \begin{enumerate}
                        \item
                            \(
                                \forall a \in A \ a \leq C
                            \)
                        \item
                            \(
                                \forall \varepsilon > 0 \ \exists a \in A : a > C - \varepsilon
                            \)
                    \end{enumerate}

                    Если \( A \) ограничено сверху, то \( \exists \sup A \)
                \item
                    Если существует \( \exists c \in \RR \ \forall a \in A \ a \geq c \),
                    то \( c \) называется нижней гранью множества \( A \).

                    Наибольшая из нижних граней множества \( A \) называется точной нижней гранью
                    и обозначается \( \inf \) (инфимум).

                    \( c = \inf A \), если:
                    \begin{enumerate}
                        \item
                            \(
                                \forall a \in A \ a \geq c
                            \)
                        \item
                            \(
                                \forall \varepsilon > 0 \ \exists a \in A : a < c + \varepsilon
                            \)
                    \end{enumerate}

                    Если \( A \) ограничено снизу, то \( \exists \inf A \)
            \end{enumerate}
    \end{itemize}

    \newpage

    \silentsubsection{6}

    \textbf{Условие:}

    Определение монотонной последовательности.
    Теорема Вейерштрасса для последовательностей.
    Определение числа \( e \).
    Определение того, что \( \lim\limits_{n \to \infty} a_n = +\infty \).
    Теорема о сходимости к \( e \) последовательностей более общего вида.

    \bigskip \bigskip

    \begin{itemize}
        \item
            Монотонная последовательность.

            Если \( \forall n \in N \ a_n \leq a_{n + 1} \), то
            \( \{ a_n \} \) называется неубывающей, если \( a_n < a_{n + 1} \) --- возрастающей.

            Если \( \forall n \in N \ a_n \geq a_{n + 1} \), то
            \( \{ a_n \} \) называется невозрастающей, если \( a_n > a_{n + 1} \) --- убывающей.
        \item
            Теорема Вейерштрасса.

            Если \( \{ a_n \} \) монотонна и ограничена, то у нее есть предел.
        \item
            \( \lim\limits_{n \to \infty} \left( 1 + \frac{1}{n} \right)^n = e \)

            \( \lim\limits_{n \to \infty} \left( 1 + \frac{1}{n} \right)^{n + 1} = e \)

            (первое стремится снизу, второе --- сверху)
        \item
            \(
                \lim\limits_{n \to \infty} a_n = +\infty
                \Leftrightarrow
                \forall M \ \exists N : \forall n > N \ a_n > M
            \)
        \item
            \begin{enumerate}
                \item
                    \( \{ p_n \} \) --- последовательность вещественных чисел,
                    и \( \lim\limits_{n \to \infty} p_n = +\infty \).

                    Тогда \( \lim\limits_{n \to \infty} \left( 1 + \frac{1}{p_n} \right)^{p_n} = e \)
                \item
                    \( \{ q_n \} \) --- последовательность вещественных чисел,
                    и \( \lim\limits_{n \to \infty} q_n = -\infty \).

                    Тогда \( \lim\limits_{n \to \infty} \left( 1 + \frac{1}{q_n} \right)^{q_n} = e \)
            \end{enumerate}
    \end{itemize}

    \newpage

    \silentsubsection{7}

    \textbf{Условие:}

    Определение подпоследовательности и частичного предела.
    Частичные пределы последовательности, имеющей предел.
    Теорема Больцано – Вейерштрасса.

    \bigskip \bigskip

    \begin{itemize}
        \item
            Подпоследовательность.

            Пусть \( n_1 < n_2 < \ldots < n_k < \ldots \) --- возрастающая последовательность натуральных чисел.

            Пусть \( \{ a_n \} \) --- последовательность, тогда последовательность \( b_k = a_{n_k} \)
            называют подпоследовательностью \( \{ a_n \} \)
        \item
            Частичный предел.

            Число \( B \) называют частичным пределом \( \{ a_n \} \),
            если существует такая ее подпоследовательность \( \{ b_k \} \),
            что \( \lim\limits_{k \to \infty} b_k = B \)

            Если \( \lim\limits_{n \to \infty} a_n = A \),
            то у любой ее подпоследовательности существует предел, равный \( A \).
        \item
            Теорема Больцано – Вейерштрасса.

            Если \( \{ a_n \} \) ограничена, то у нее есть частичный предел.
    \end{itemize}

    \newpage

    \silentsubsection{8}

    \textbf{Условие:}

    Определение верхнего и нижнего предела последовательности.
    Структура множества частичных пределов.
    Критерий существования предела в терминах частичных пределов.

    \bigskip \bigskip

    \begin{itemize}
        \item
            Верхний и нижний пределы

            Пусть \( \{ a_n \} \) ограничена. \( M_n = \sup \{ a_{n + 1}, a_{n + 2}, \ldots \} \).

            Тогда \( \exists \lim\limits_{n \to \infty} M_n = M \).

            Он называется верхним пределом последовательности \( \{ a_n \} \).

            Аналогично \( m_n = \inf \{ a_{n + 1}, a_{n + 2}, \ldots \} \).

            \( \lim\limits_{n \to \infty} m_n = m \) называется нижним пределом последовательности \( \{ a_n \} \).
        \item
            Структура множества частичных пределов.

            \( m, M \) --- частичные пределы \( \{ a_n \} \).

            Если \( b \) --- частичный предел \( \{ a_n \} \), то \( m \leq b \leq M \).
        \item
            Критерий существования предела через частичные пределы.

            Ограниченная последовательность \( \{ a_n \} \) имеет предел
            \(
                \Leftrightarrow
                \underline{\lim} \ a_n = \overline{\lim} \ a_n
            \)
    \end{itemize}

    \newpage

    \silentsubsection{9}

    \textbf{Условие:}

    Определение фундаментальной последовательности.
    Формулировка критерия Коши.
    Определение числового ряда, частичной суммы, последовательности частичных сумм и суммы ряда.

    \bigskip \bigskip

    \begin{itemize}
        \item
            Фундаментальная последовательность.

            \( \{ a_n \} \) --- фундаментальная последовательность, если
            \[
                \forall \varepsilon > 0 \ \exists N : \forall n, m > N \ |a_n - a_m| < \varepsilon
            \]
        \item
            Критерий Коши

            \( \{ a_n \} \) имеет предел \( \Leftrightarrow \) \( \{ a_n \} \) --- фундаментальная.
        \item
            Числовые ряды.

            Пусть дана \( \{ a_n \} \).

            Формальная сумма \( a_1 + \ldots + a_n + \ldots = \sum\limits_{n = 1}^{\infty} a_n \)
            называется числовым рядом.

            Последовательность \( S_N = \sum\limits_{n = 1}^N a_n \)
            называется последовательностью частичных сумм ряда.

            Элементы последовательности \( \{ a_n \} \) называются элементами ряда.

            Если \( \exists S = \lim\limits_{N \to \infty} S_N \),
            то \( S \) называется суммой ряда.

            Если \( S \) существует, то говорят, что ряд сходится.

            Еще есть разные суммирования (вероятно, ненужно):
            \begin{itemize}
                \item
                    По Чезаро:
                    \[
                        S = \lim\limits_{n \to \infty} \frac{S_1 + \ldots + S_n}{n}
                    \]
                \item
                    По Пуассону-Абелю:
                    \[
                        S = \lim\limits_{x \to 1-} \sum\limits_{n = 1}^{\infty} a_n x^n
                    \]
            \end{itemize}
    \end{itemize}

    \newpage

    \silentsubsection{10}

    \textbf{Условие:}

    Критерий Коши сходимости ряда.
    Необходимый признак сходимости ряда.
    Абсолютная и условная сходимость ряда.
    Критерий сходимости рядов с неотрицательными членами

    \bigskip \bigskip

    \begin{itemize}
        \item
            Критерий Коши сходимости ряда.

            Ряд сходится
            \(
                \Leftrightarrow
                \forall \varepsilon > 0 \ \exists N : \forall m > n > N
                \ |S_m - S_n| = \left| \sum\limits_{k = n + 1}^m a_k \right| < \varepsilon
            \)
        \item
            Необходимый признак сходимости ряда.

            Ряд \( \sum\limits_{n = 1}^{\infty} a_n \) сходится \( \Rightarrow \lim\limits_{n \to \infty} a_n = 0 \)
        \item
            Абсолютная и условная сходимость ряда.

            Пусть ряд \( \sum\limits_{n = 1}^{\infty} a_n \) сходится.

            Если ряд \( \sum\limits_{n = 1}^{\infty} |a_n| \) тоже сходится, то такую сходимость называют абсолютной,
            в противном случае --- условной.
        \item
            Критерий сходимости рядов с неотрицательными членами.

            Ряд сходится \( \Leftrightarrow \) \( \{ S_n \} \) --- ограниченная последовательность.
    \end{itemize}

    \newpage

    \silentsubsection{11}

    \textbf{Условие:}

    Признак сравнения.
    Признак сравнения в предельной форме.
    Мажорантный признак Вейерштрасса.
    Признак разрежения Коши.
    Признак Даламбера.
    Признак Коши.
    Признак Гаусса.

    \bigskip \bigskip

    \begin{itemize}
        \item
            Признак сравнения.

            Пусть \( \exists n_0 \in \NN : \forall n > n_0 \ 0 \leq a_n \leq b_n \),

            \( \sum\limits_{n = 1}^{\infty} b_n \) сходится \( \Rightarrow \) \( \sum\limits_{n = 1}^{\infty} a_n \) сходится.

            \( \sum\limits_{n = 1}^{\infty} a_n \) расходится \( \Rightarrow \) \( \sum\limits_{n = 1}^{\infty} b_n \) расходится.
        \item
            Признак сравнения в предельной форме.

            Пусть \( \exists N : \forall n > N \ a_n \geq 0, b_n > 0 \)
            и \( \lim\limits_{n \to \infty} \frac{a_n}{b_n} = A \in (0, +\infty) \)

            Тогда либо \( \sum\limits_{n = 1}^{\infty} a_n, \ \sum\limits_{n = 1}^{\infty} b_n \) оба сходятся,
            либо оба расходятся.
        \item
            Мажорантный признак Вейерштрасса.

            Пусть \( \forall n > N \ |a_n| \leq b_n \)

            \( \sum\limits_{n = 1}^{\infty} b_n \) сходится \( \Rightarrow \sum\limits_{n = 1}^{\infty} a_n \) тоже сходится.
        \item
            Признак разрежения Коши.

            Пусть \( \{ a_n \} \) не возрастает.

            Тогда \( \sum\limits_{n = 1}^{\infty} a_n \) сходится тогда и только тогда,
            когда сходится \( \sum\limits_{k = 1}^{\infty} 2^k a_{2^k} \)
        \item
            Признак Д'Аламбера.

            Пусть дан ряд \( \sum\limits_{n = 1}^{\infty} a_n \)
            и \( \exists \lim\limits_{n \to \infty} \left| \frac{a_{n + 1}}{a_n} \right| = q \).
            \begin{enumerate}
                \item
                    \( q < 1 \Rightarrow \sum\limits_{n = 1}^{\infty} a_n \) сходится
                \item
                    \( q > 1 \Rightarrow \sum\limits_{n = 1}^{\infty} a_n \) расходится
                \item
                    \( \not\exists\lim \) или \( q = 1 \Rightarrow \) может быть что угодно
            \end{enumerate}

            \newpage
        \item
            Радикальный Признак Коши.

            Пусть дан ряд \( \sum\limits_{n = 1}^{\infty} a_n \)
            и \( \exists \overline{\lim\limits_{n \to \infty}} \sqrt[n]{|a_n|} = q \)

            Тогда
            \begin{enumerate}
                \item
                    \( q < 1 \Rightarrow \sum\limits_{n = 1}^{\infty} a_n \) сходится
                \item
                    \( q > 1 \Rightarrow \sum\limits_{n = 1}^{\infty} a_n \) расходится
                \item
                    \( \not\exists\lim \) или \( q = 1 \Rightarrow \) может быть что угодно
            \end{enumerate}
        \item
            Признак Гаусса.

            Пусть дан ряд \( \sum\limits_{n = 1}^{\infty} a_n \) и \( a_n \geq 0 \ \forall n \in \NN \)

            Пусть
            \[
                \exists \delta > 0 : \frac{a_n}{a_{n + 1}} = \alpha + \frac{\beta}{n} + \gamma_n \cdot \frac{1}{n^{1 + \delta}}
            \]

            причем \( \gamma_n \) --- ограниченная последовательность.

            \begin{enumerate}
                \item
                    \( \alpha < 1 \Rightarrow \) расходится
                \item
                    \( \alpha > 1 \Rightarrow \) сходится
                \item
                    \( \alpha = 1 \)

                    \begin{enumerate}
                        \item
                            \( \beta > 1 \Rightarrow \) сходится
                        \item
                            \( \beta \leq 1 \Rightarrow \) расходится
                    \end{enumerate}
            \end{enumerate}
    \end{itemize}

    \newpage

    \silentsubsection{12}

    \textbf{Условие:}

    Определение покрытия.
    Принцип Бореля-Лебега.
    Предельная точка множества (оба определения).
    Принцип Больцано-Вейерштрасса.
    Внутренняя, граничная и изолированная точка.
    Открытое и замкнутое множество.
    Критерий замкнутости множества.
    Определение компакта.

    \bigskip \bigskip

    \begin{itemize}
        \item
            Покрытие

            Система множеств \( S = \{ U_\alpha \} \) называется покрытием множества \( A \), \\
            если \( A \subseteq \bigcup\limits_{U_\alpha \in S} U_\alpha \)

        \item
            Принцип Бореля-Лебега

            Из любого покрытия отрезка \( [a, b] \) системы интервалов \( S = \{ J_\alpha \} \)
            можно выбрать подсистему в \( S \), состоящую из конечного числа интервалов и покрывающую \( [a, b] \).
        \item
            Предельная точка

            Пусть дано множество \( A \).

            Тогда \( a \) называется предельной точной множества \( A \),
            если \( \forall \mathring{U}(a) \ \exists b \in A : b \in \mathring{U}(a) \)

            Тогда \( a \) называется предельной точной множества \( A \),
            если \( \forall U(a) \) существует бесконечно много элементов из \( A \), лежащих в \( U(a) \)
        \item
            Принцип Больцано-Вейерштрасса

            Пусть \( A \) --- бесконечное множество (бесконечное число элементов),
            причем ограниченное (то есть \( \exists C > 0 : \forall a \in A \ |a| \leq C \)

            Тогда существует хотя бы одна предельная точка множества \( A \)
        \item
            Внутренняя, граничная и изолированная точка.

            Если \( a \in A \) и \( \exists U(a) : U(a) \subseteq A \), то \( a \) --- внутренняя точка множества \( A \).

            Если \( \forall U(b) \ \exists \) точки из \( A \) и точки из \( \RR \setminus A \),
            лежащие в \( U(b) \), то \( b \) называется граничной точкой множества \( A \).

            Если \( c \in A \) и \( \exists U(c) : U(c) \cap A = c \), то \( c \) называется изолированной.
        \item
            Открытое и замкнутое множество.

            Множество открыто, если любая его точка является внутренней.

            Множество \( A \) замкнуто, если \( \RR \setminus A \) открыто.
        \item
            Критерий замкнутости множества.

            \( A \) замкнуто \( \Leftrightarrow \) \( A \) содержит все свои предельные точки.
        \item
            Компакт.

            Множество \( K \subset \RR \) называется компактом,
            если из любого его покрытия можно выбрать конечное подпокрытие.
    \end{itemize}

    \newpage

    \silentsubsection{13}

    \textbf{Условие:}

    Определение предела по Коши (в том числе через окрестности).
    Определение предела при \( x \to +\infty \).
    Определение предела по Гейне.
    Эквивалентность определений по Коши и Гейне.

    \bigskip \bigskip

    \begin{itemize}
        \item
            Определение предела по Коши.

            Пусть \( f : E \to \RR \).

            Пусть \( a \in E \).
            \begin{gather*}
                \lim\limits_{x (\in E) \to a} f(x) = A
                \\
                \Updownarrow
                \\
                \forall \varepsilon > 0 \ \exists \delta > 0 : \forall x :
                    (x \in E \land 0 < |x - a| < \delta) \ |f(x) - A| < \varepsilon
            \end{gather*}

            Через окрестности:
            \[
                \forall V_{\varepsilon}(A) \ \exists \mathring{U_\delta} (a) :
                \forall x \in E \cap \mathring{U_\delta} (a) \ f(x) \in V_{\varepsilon} (A)
            \]
        \item
            Определение предела при \( x \to +\infty \).
            \[
                \forall U_{\varepsilon}(A) \ \exists M :
                \forall x \in E \cap (M, +\infty) \ f(x) \in U_{\varepsilon} (A)
            \]
        \item
            Определение предела по Гейне.
            \begin{gather*}
                \lim\limits_{x (\in E) \to a} f(x) = A
                \\
                \Updownarrow
                \\
                \forall \{ a_n \} \in E : ( \ \lim\limits_{n \to \infty} a_n = a \land a_n \neq a \ \forall n \in \NN )
                \ \lim\limits_{n \to \infty} f(a_n) = A
            \end{gather*}
        \item
            Эквивалентность определений по Коши и Гейне.

            Два вышеприведенных определения эквивалентны.
    \end{itemize}

    \newpage

    \silentsubsection{14}

    \textbf{Условие:}

    Свойства предела.
    Теорема о пределе композиции.
    Эквивалентные функции.

    \bigskip \bigskip

    \begin{itemize}
        \item
            Свойства предела.

            Пусть \( f, g, h : E \to \RR \).

            \( a \in E, \lim\limits_{x \to a} f(x) = A, \ \lim\limits_{x \to b} g(x) = B \).

            \begin{enumerate}
                \item
                    Единственность
                \item
                    \( \forall \alpha, \beta \in \RR \ \lim\limits_{x \to a} (\alpha f(x) + \beta g(x)) = \alpha A + \beta B \)
                \item
                    \( \lim\limits_{x \to a} f(x)g(x) = AB \)
                \item
                    Пусть
                    \(
                        B \neq 0, \ g(x) \neq 0 \ \forall x \in E \Rightarrow \lim\limits_{x \to a} \frac{f(x)}{g(x)} = \frac{A}{B}
                    \)
                \item
                    Предельный переход в неравенстве.

                    (нужна окрестность с \( f(x) \leq g(x) \))
                \item
                    Если функция имеет предел в точке, то существует окрестность этой точки, в которой функция ограничена.
                \item
                    Лемма о зажатом пределе

                    (нужна окрестность с \( f(x) \leq h(x) \leq g(x)\) и \( A = B \))
                \item
                    Отделимость от нуля (можем найти ``хорошую'' окрестность).
            \end{enumerate}
        \item
            Теорема о пределе композиции.
            \begin{gather*}
                g : D \to \RR, \ b \in D
                \\
                f : E \to D, \ a \in E
                \\
                \exists \mathring{U}(a) : \forall x \in \mathring{U}(a) \ f(x) \neq b
                \\
                \lim\limits_{x \to a} f(x) = b, \ \lim\limits_{y \to b} g(y) = c \Rightarrow \lim\limits_{x \to a} g(f(x)) = c
            \end{gather*}
        \item
            Эквивалентные функции.

            Пусть \( f, g : E \to \RR \) и \( a \) --- предельная точка \( E \)

            \( f \sim g, \ x \to a \), если
            \begin{enumerate}
                \item
                    \( \exists \mathring{U}(a) : \forall x \in \mathring{U}(a) \cap E \ f(x) = g(x)h(x) \),
                    где \( h : \mathring{U}(a) \cap E \to \RR \)
                \item
                    \( \lim\limits_{x \to a} h(x) = 1 \)
            \end{enumerate}

            При \( g(x) \neq 0 \) в \( \mathring{U}(a) \cap E \) это равносильно
            \( \lim\limits_{x \to a} \frac{f(x)}{g(x)} = 1 \)
    \end{itemize}

    \newpage

    \silentsubsection{15}

    \textbf{Условие:}

    Критерий Коши существования предела функции.
    Определение односторонних пределов.
    Определение монотонной функции.
    Определение ограниченной функции.
    Теорема Вейерштрасса для монотонной функции.

    \bigskip \bigskip

    \begin{itemize}
        \item
            Критерий Коши существования предела функции.
            \[
                \exists \lim\limits_{x \to a} f(x)
                \Leftrightarrow
                \forall \varepsilon > 0 \ \exists \delta > 0 :
                    \forall x_1, x_2 \in \mathring{U}_\delta(a) \ |f(x_1) - f(x_2)| < \varepsilon
            \]
        \item
            Односторонние пределы.
            \begin{gather*}
                E_a^- = \{ x \in E : x < a \}
                \\
                E_a^+ = \{ x \in E : x > a \}
            \end{gather*}

            \( \lim\limits_{x \to a-} f(x) = \lim\limits_{x ( \in E_a^- ) \to a} \) --- предел слева.

            \( \lim\limits_{x \to a+} f(x) = \lim\limits_{x ( \in E_a^+ ) \to a} \) --- предел справа.
        \item
            Монотонные функции.

            Пусть \( f : E \to \RR \) и \( \forall x_1, x_2 \in E : x_1 < x_2 \)
            \begin{enumerate}
                \item
                    \( f(x_1) \leq f(x_2) \Rightarrow f \uparrow \)
                \item
                    \( f(x_1) < f(x_2) \Rightarrow f \uparrow \uparrow \)
                \item
                    \( f(x_1) \geq f(x_2) \Rightarrow f \downarrow \)
                \item
                    \( f(x_1) > f(x_2) \Rightarrow f \downarrow \downarrow \)
            \end{enumerate}
        \item
            Ограниченная функции.

            Функция \( f \) ограничена на множестве \( A \),
            если \( f \) определена на \( A \),
            и существует константа \( C : \forall x \in A \ |f(x)| < C \).
        \item
            Теорема Вейерштрасса для монотонной функции.

            Пусть \( f : E \to \RR \)
            \begin{enumerate}
                \item
                    Если \( f \uparrow \) и ограничена сверху на \( E_a^- \),
                    а \( a \) --- предельная точка \( E_a^- \),
                    то \( \exists \lim\limits_{x \to a-} f(x) \)
                \item
                    Если \( f \downarrow \) и ограничена снизу \( E_a^+ \),
                    а \( a \) --- предельная точка \( E_a^+ \),
                    то \( \exists \lim\limits_{x \to a+} f(x) \)
            \end{enumerate}
    \end{itemize}

    \newpage

    \silentsubsection{16}

    \textbf{Условие:}

    Определение бесконечно малой функции.
    Определение \( o \)-малого.
    \( o \)-малое в терминах пределов.
    Определение \( O \)-большого.
    Асимптотические равенства \( 1 \) - \( 13 \).

    \bigskip \bigskip

    \begin{itemize}
        \item
            Бесконечно малая функция.

            \( f : E \to \RR \) называется бесконечно малой при \( x \to a \),
            если \( \lim\limits_{x (\in E) \to a} f(x) = 0 \)
        \item
            \( o \)-малое и \( O \)-большое

            Пусть \( f, g : E \to \RR \) и \( a \) --- предельная точка \( E \).

            Пусть \( \exists \mathring{U}(a) : \forall x \in \mathring{U}(a) \cap E \ f(x) = g(x)h(x) \).

            Тогда:
            \begin{enumerate}
                \item
                    Если \( \lim\limits_{x \to a} h(x) = 0 \), то \( f = o(g), x \to a \)
                \item
                    Если \( h \) ограничено на \( \mathring{U}(a) \cap E \), то \( f = O(g), x \to a \)
            \end{enumerate}

            Если в \( \mathring{U}(a) \cap E \ g(x) \neq 0 \), то
            \begin{enumerate}
                \item
                    \( \Leftrightarrow \lim\limits_{x \to a} \frac{f(x)}{g(x)} = 0 \)
                \item
                    \( \Leftrightarrow \frac{f(x)}{g(x)} \) --- ограниченная функция в \( \mathring{U}(a) \cap E \)
            \end{enumerate}

            Тогда бесконечно малая при \( x \to a \) может быть обозначена как \( o(1) \).
        \item
            Асимптотические равенства.

            \begin{enumerate}
                \item
                    \( \sin x = x + o(x), \ x \to 0 \)
                \item
                    \( \tan x = x + o(x), \ x \to 0 \)
                \item
                    \( \arctan x = x + o(x), \ x \to 0 \)
                \item
                    \( \arcsin x = x + o(x), \ x \to 0 \)
                \item
                    \( \ln (1 + x) = x + o(x), \ x \to 0 \)
                \item
                    \( e^x = 1 + x + o(x), \ x \to 0 \)
                \item
                    \( \cos x = 1 - \frac{x^2}{2} + o(x), \ x \to 0 \)
                \item
                    \( (1 + x)^\alpha = 1 + \alpha x + o(x), \ x \to 0 \)
            \end{enumerate}

            \newpage

            Обрубки ряда Тейлора в \( x = 0 \):
            \begin{enumerate}
                \item
                    \(
                        \displaystyle
                        e^x = \sum\limits_{k = 0}^n \frac{x^k}{k!} + o(x^n)
                    \)
                \item
                    \(
                        \displaystyle
                        \sin x = \sum\limits_{k = 0}^n (-1)^k \frac{x^{2k + 1}}{(2k + 1)!} + o(x^{2n + 1})
                    \)
                \item
                    \(
                        \displaystyle
                        \cos x = \sum\limits_{k = 0}^n (-1)^k \frac{x^{2k}}{(2k)!} + o(x^{2n})
                    \)
                \item
                    \(
                        \displaystyle
                        \ln (1 + x) = \sum\limits_{k = 1}^n (-1)^k \frac{x^k}{k} + o(x^n)
                    \)
                \item
                    \(
                        \displaystyle
                        (1 + x)^\alpha = \sum\limits_{k = 0}^n \frac{\alpha (\alpha - 1) \ldots (\alpha - k + 1)}{k!} x^k + o(x^n)
                    \)
            \end{enumerate}
    \end{itemize}

    \newpage

%%%%%%%%%%%%%%%%%%%%%%%%%%%%%%%%%%%%%%%%%%%%%%%%%%%%%%%%%%%%%%%%%%%%%%%%%%%%%%%%%%%%%%%%%%%%%%%%%%%%%%%%%%%%%%%%%%%%%%%%%%%

    \silentsection{Листик 1}

    \begin{center}
        \textit{(Индукция, суммирование, действительные числа)}
    \end{center}

    \silentsubsection{Вопрос 1}

    \textbf{Условие:}

    Доказать равенства и неравенства по индукции:

    \silentsubsubsection{а)}

    \textbf{Условие:}

    \( 1^3 + 2^3 + \dots + n^3 = \left( \frac{n(n + 1)}{2} \right)^2\)

    \textit{(0.5 балла)}

    \bigskip \bigskip

    Докажем по индукции:

    База: \( n = 1, \ 1^3 = \left( \frac{1 \cdot 2}{2} \right)^2\)

    Переход:

    Пусть верно для \( n = k \), докажем для \( n = k + 1 \)
    \[ 1^3 + 2^3 + \dots + k^3 = \left( \frac{k(k + 1)}{2} \right)^2 \]

    Добавим \( (k + 1)^3 \) к обеим сторонам равенства:
    \begin{gather*}
        1^3 + 2^3 + \dots + k^3 + (k + 1)^3 = \left( \frac{k(k + 1)}{2} \right)^2 + (k + 1)^3 = \\
        = \frac{k^2 (k + 1)^2 + 4(k + 1)^3}{4} = \frac{(k + 1)^2 (k^2 + 4k + 4)}{4} = \\
        \frac{(k + 1)^2 (k + 2)^2}{4} = \left( \frac{(k + 1)(k + 2)}{2} \right)^2
    \end{gather*}

    Что и требовалось доказать.
    \newpage

    \silentsubsubsection{б)}

    \textbf{Условие:}

    \( (x + y)^n = \sum\limits_{i = 0}^n C_n^i x^i y^{n - i} \)

    \textit{(0.5 балла)}

    \bigskip \bigskip

    Докажем по индукции:

    База: \(n = 1, \ (x + y)^1 = C_1^0 x^0 y^1 + C_1^1 x^1 y^0 \)

    Переход:

    Пусть верно для \( n = k \), докажем для \( n = k + 1 \)
    \[ (x + y)^k = \sum\limits_{i = 0}^k C_k^i x^i y^{k - i} \]

    Домножим на \( (x + y) \)

    \begin{gather*}
        (x + y)^{k + 1} = (x + y) \sum\limits_{i = 0}^k C_k^i x^i y^{k - i} = \\
        = \sum\limits_{i = 0}^k C_k^i (x^{i + 1} y^{k - i} + x^i y^{k + 1 - i})
            = C_k^0 y^{k + 1} + \sum\limits_{i = 1}^k (C_k^{i - 1} + C_k^i) x^i y^{k + 1 - i} + C_k^k x^{k + 1} = \\
        = C_{k + 1}^0 y^{k + 1} + \sum\limits_{i = 1}^k C_{k + 1}^i x^i y^{k + 1 - i}
            + C_{k + 1}^{k + 1} x^{k + 1} = \\
        = \sum\limits_{i = 0}^{k + 1} C_{k + 1}^i x^i y^{k + 1 - i}
    \end{gather*}

    Что и требовалось доказать.
    \newpage

    \silentsubsubsection{в)}

    \textbf{Условие:}

    \( (n + 1)^n < n^{n + 1}, \ n \geq 3 \)

    \textit{(0.5 балла)}

    \bigskip \bigskip

    Докажем по индукции:

    База: \( n = 3, 4^3 = 64 < 81 = 3^4 \)

    Переход:

    Пусть верно для \( n = k \), докажем для \( n = k + 1 \)

    \[ (k + 1)^k < k^{k + 1} \]

    Домножим на \( \frac{(k + 2)^{k + 1}}{(k + 1)^k} > 0 \)

    \[ (k + 2)^{k + 1} < k^{k + 1} \cdot \frac{(k + 2)^{k + 1}}{(k + 1)^k} \]

    Покажем, что
    \[ k^{k + 1} \cdot \frac{(k + 2)^{k + 1}}{(k + 1)^k} < (k + 1)^{k + 2} \]

    Действительно:
    \begin{gather*}
        k^{k + 1} \cdot \frac{(k + 2)^{k + 1}}{(k + 1)^k} < (k + 1)^{k + 2} \\
        k^{k + 1} (k + 2)^{k + 1} < (k + 1)^{2k + 2} \\
        k(k + 2) < (k + 1)^2 \\
        k^2 + 2k < k^2 + 2k + 1 \\
        0 < 1
    \end{gather*}

    Значит
    \[
        (k + 2)^{k + 1} < k^{k + 1} \cdot \frac{(k + 2)^{k + 1}}{(k + 1)^k} < (k + 1)^{k + 2}
    \]

    Что и требовалось доказать.
    \newpage

    \silentsubsubsection{г)}

    \textbf{Условие:}

    \( \sqrt{n} < \sum\limits_{i = 1}^n \frac{1}{\sqrt{i}} < 2\sqrt{n}, \ n > 1 \)

    \textit{(0.5 балла)}

    \bigskip \bigskip

    Докажем по индукции:

    База:
    \(
        n = 2, \ \sqrt{2} = 2 \cdot \frac{1}{\sqrt{2}} < \frac{1}{\sqrt{1}} + \frac{1}{\sqrt{2}}
        < 2 \cdot 1 < 2 \sqrt{2}
    \)

    Переход:

    Пусть верно для \( n = k \), докажем для \( n = k + 1 \)

    \[ \sqrt{k} < \sum\limits_{i = 1}^k \frac{1}{\sqrt{i}} < 2\sqrt{k} \]

    Прибавим \( \frac{1}{\sqrt{k + 1}} \)
    \[
        \sqrt{k} + \frac{1}{\sqrt{k + 1}} < \sum\limits_{i = 1}^{k + 1} \frac{1}{\sqrt{i}}
        < 2 \sqrt{k} + \frac{1}{\sqrt{k + 1}}
    \]

    Докажем, что
    \[
        \sqrt{k + 1} < \sqrt{k} + \frac{1}{\sqrt{k + 1}}
        \quad \text{и} \quad
        2 \sqrt{k} + \frac{1}{\sqrt{k + 1}} < 2 \sqrt{k + 1}
    \]

    Действительно:

    \begin{gather*}
        \sqrt{k + 1} < \sqrt{k} + \frac{1}{\sqrt{k + 1}} \\
        k + 1 < \sqrt{k (k + 1)} + 1 \\
        k < \sqrt{k (k + 1)} \\
        k^2 < k^2 + k \\
        0 < k
    \end{gather*}

    \begin{gather*}
        2 \sqrt{k} + \frac{1}{\sqrt{k + 1}} < 2 \sqrt{k + 1} \\
        2 \sqrt{k (k + 1)} + 1 < 2 (k + 1) \\
        2 \sqrt{k (k + 1)} < 2k + 1 \\
        4 (k^2 + k) < 4k^2 + 4k + 1 \\
        4k^2 + 4k < 4k^2 + 4k + 1 \\
        0 < 1
    \end{gather*}

    Тогда

    \[
        \sqrt{k + 1}
        < \sqrt{k} + \frac{1}{\sqrt{k + 1}}
        < \sum\limits_{i = 1}^{k + 1} \frac{1}{\sqrt{i}}
        < 2 \sqrt{k} + \frac{1}{\sqrt{k + 1}}
        < 2 \sqrt{k + 1}
    \]

    Что и требовалось доказать
    \newpage

    \silentsubsection{Вопрос 2}

    \textbf{Условие:}

    Найти суммы:

    \silentsubsubsection{а)}

    \textbf{Условие:}

    \( C_{100}^0 + C_{100}^1 + \dots + C_{100}^{100} \)

    \textit{(0.5 балла)}

    \bigskip \bigskip

    Решим комбинаторно:

    \( C_{100}^i \) --- количество подмножеств набора из \( 100 \) элементов размера \( i \).
    Так как суммируем по всем \( i \) от \( 0 \) до \( 100 \),
    ответом будет количество всех подмножеств набора из \( 100 \) элементов --- \( 2^{100} \)

    \bigskip \bigskip

    Решим через бином:
    \[
        2^{100} = (1 + 1)^{100} = \sum\limits_{k = 0}^{100} 1^{100} C_{100}^k
        =
        C_{100}^0 + C_{100}^1 + \dots + C_{100}^{100}
    \]

    \textbf{Ответ:} \( 2^{100} \)
    \newpage

    \silentsubsubsection{б)}

    \textbf{Условие:}

    \( \sum\limits_{k = 1}^n 2^k C_n^k \)

    \textit{(0.5 балла)}

    \bigskip \bigskip

    \( \sum\limits_{k = 1}^n 2^k C_n^k = \sum\limits_{k = 0}^n 2^k C_n^k - 1\)

    Комбинаторно найдем \( \sum\limits_{k = 0}^n 2^k C_n^k \)

    \( C_n^k \) --- выбрали подмножество размера \( k \).
    \( 2^k \) --- покрасили его в два цвета.
    Оставшееся покрасим в третий цвет.
    Таким образом, одно слагаемое суммы
    --- количество способов покрасить \( n \) элементов в \( 3 \) цвета так,
    чтобы элементов первого и второго цвета было \( k \).
    Тогда вся сумма --- количество способов покрасить \( n \) элементов в \( 3 \) цвета,
    т.е. \( \sum\limits_{k = 0}^n 2^k C_n^k = 3^n \). Значит ответ на задачу --- \( 3^n - 1\).

    \bigskip \bigskip

    Решим через бином:
    \[
        3^n = (2 + 1)^n = \sum\limits_{k = 0}^n 1^{n-k} 2^k C_n^k
        =
        \sum\limits_{k = 1}^n 2^k C_n^k + 1
    \]

    \textbf{Ответ:} \(3^n - 1\)
    \newpage

    \silentsubsubsection{в)}

    \textbf{Условие:}

    \( C_n^0 - C_n^1 + C_n^2 - \dots + (-1)^n C_n^n \)

    \textit{(0.5 балла)}

    \bigskip \bigskip

    Заметим, что ответ равен разнице количеств подмножеств чётного размера
    и подмножеств нечётного размера набора из \( n \) элементов.

    \( n \) --- натуральное, а значит набор не пустой.

    Разобьем подмножества на пары: каждому подмножеству, содержащему \( 1 \)-й элемент, сопоставим
    такой же набор, но без \( 1 \)-го элемента.
    Очевидно, что таким образом мы разобьем абсолютно все подмножества на пары
    (биекцию несложно проверить), причем в каждой паре размеры подмножеств различаются на \( 1 \).
    Это означает, что количества подмножеств с чётным размером и с нечётным размером равны
    \( \Rightarrow \) сумма из условия равна \( 0 \).

    \bigskip \bigskip

    Решим через бином (\( n > 0\)):

    \[
        0 = (1 + (-1))^n = \sum\limits_{k = 0}^n 1^{n-k} (-1)^k C_n^k
        =
        C_n^0 - C_n^1 + C_n^2 - \dots + (-1)^n C_n^n
    \]

    \textbf{Ответ:} \( 0 \) (если вдруг \( n = 0\), то \( 1 \))
    \newpage

    \silentsubsection{Вопрос 3}

    \textbf{Условие:}

    Доказать, что \( (\sqrt{3} - \sqrt{2})^{999} \) можно представить в виде \( a \sqrt{3} - b \sqrt{2}\),
    причем \( 3a^2 - 2b^2 = 1\).

    \textit{(1.5 балла)}

    \bigskip \bigskip

    \(
        (\sqrt{3} - \sqrt{2}) (\sqrt{3} + \sqrt{2}) = 1 \Rightarrow \sqrt{3} + \sqrt{2} =
        \frac{1}{\sqrt{3} - \sqrt{2}}
    \)

    \begin{gather*}
        3a^2 - 2b^2 = 1 \\
        (a \sqrt{3} - b \sqrt{2}) (a \sqrt{3} + b \sqrt{2}) = 3 - 2 = 1
    \end{gather*}

    Если \( (\sqrt{3} - \sqrt{2})^{999} = a \sqrt{3} - b \sqrt{2} \), то
    \begin{gather*}
        (\sqrt{3} - \sqrt{2})^{999}(a \sqrt{3} + b \sqrt{2}) = 1 \\
        a \sqrt{3} + b \sqrt{2} = \frac{1}{(\sqrt{3} - \sqrt{2})^{999}} = (\sqrt{3} + \sqrt{2})^{999}
    \end{gather*}

    Тогда если мы решим

    \[
        \begin{cases}
            a \sqrt{3} - b \sqrt{2} = (\sqrt{3} - \sqrt{2})^{999} \\
            a \sqrt{3} + b \sqrt{2} = (\sqrt{3} + \sqrt{2})^{999}
        \end{cases}
    \]

    то очевидно, что такие \( a, b \) подойдут под условие.

    Решением является:
    \[
        \begin{cases}
            a = \frac{(\sqrt{3} - \sqrt{2})^{999} + (\sqrt{3} + \sqrt{2})^{999}}{2 \sqrt{3}} \\
            b = \frac{(\sqrt{3} + \sqrt{2})^{999} - (\sqrt{3} - \sqrt{2})^{999}}{2 \sqrt{2}}
        \end{cases}
    \]

    А значит подходящие \( a, b \) существуют. Что и требовалось доказать.

    \bigskip

    P.S. Если \( 999 \) заменить на \( n \), ничего не поменяется.

    \newpage

    \silentsubsection{Вопрос 4}

    \textbf{Условие:}

    Для любого ли конечного множества действительных чисел \( a_1, a_2, \ldots, a_n\)
    существует такое \(m \in  \{0, 1, 2, \ldots, n\} \), что
    \[
        \left| \sum\limits_{k = 1}^{m} a_k - \sum\limits_{k = m + 1}^{n} a_k \right|
        \leq
        \max\limits_{1 \leq l \leq n} | a_l |
    \]
    (при \( m = 0 \) считаем равной нулю первую сумму, а при \( m = n \) --- вторую)?

    \textit{(1.5 балла)}

    \bigskip \bigskip

    Обозначим левую часть неравенства без модуля за \( f(m) \), а правую --- за \( C \).
    \begin{itemize}
        \item
            \( | f(0) | \leq C \Rightarrow \) подходящее \( m \) существует.
        \item
            \( | f(0) | > C \)

            Не умаляя общности, \( f(0) > C \).

            \( f(n) = -f(0) \Rightarrow f(n) < -C \).

            Предположим, что нужного \( m \) не существует.
            В таком случае
            \begin{gather*}
                \exists x \in \{ 0, 1, \ldots, n - 1 \} : ( f(x) > C ) \land ( f(x + 1) < -C )
                \\
                \begin{cases}
                    f(x) > C
                    \\
                    f(x + 1) < -C
                \end{cases}
                \Rightarrow
                f(x) - f(x + 1) > 2C
                \Rightarrow
                -2 a_{x + 1} > 2C
                \Rightarrow
                a_{x + 1} < -C
            \end{gather*}

            Но тогда \( | a_{x + 1} | > C = \max\limits_{0 \leq l \leq n} | a_l | \) --- противоречие.

            Значит нужное \( m \) существует.
    \end{itemize}

    \textbf{Ответ:} да, всегда

    \newpage

    \silentsection{Листик 2}

    \begin{center}
        \textit{(Действительные числа)}
    \end{center}

    \silentsubsection{Вопрос 1}

    \textbf{Условие:}

    Найти все положительные числа, представимые двумя способами в виде бесконечной дроби с основанием \( q \)

    \textit{(1 балл)}

    \bigskip \bigskip

    Как мы строим запись неотрицательного числа?

    \begin{itemize}
        \item
            Разбиваем \( [0, +\infty) \) на отрезки длины \( 1 \): \( [0, 1], [1, 2], \ldots \)
        \item
            Целая часть --- номер отрезка (в \( 0 \)-индексации)
        \item
            Для получения дробной части:

            Бьем наш отрезок, в который попало число, еще на \( q \) равных, первая цифра дробной части
            --- снова номер отрезка (в \( 0 \)-индексации).
            Чтобы получить вторую часть, снова бьем отрезок, в который теперь попало число на \( q \) равных,
            определяем вторую цифру дробной части, повторяем процедуру до бесконечности.
    \end{itemize}

    Когда будет неоднозначность?

    В тот момент, когда число попадет на стык двух отрезков, после чего мы не сможем однозначно выбрать отрезок,
    в котором оно лежит.
    В таком случае число представимо двумя способами --- бесконечной дробью с периодом \( (0) \)
    или с периодом \( (q - 1) \).

    Нетрудно понять, что все такие положительные числа представимы в виде
    \[
        \frac{n}{q^k}, n \in \NN, k \in \NN_0
    \]

    Точная формулировка ответа будет немного нуднее, так как при определении рациональных чисел мы считали,
    что числитель и знаменатель взаимнопросты.

    \textbf{Ответ:}

    Пусть \( p_1, p_2, \ldots, p_k \) --- все различные простые делители \( q \).

    Число \( x \) представимо двумя способами в виде бесконечной дроби с основанием \( q \)
    тогда и только тогда, когда
    \begin{gather*}
        x = \frac{n}{p_1^{d_1} p_2^{d_2} \ldots p_k^{d_k}}
        \\
        n \in \NN, p_1, p_2, \ldots, p_k \in \NN_0
        \\
        (x, p_1) = (x, p_2) = \ldots = (x, p_k) = 1
    \end{gather*}

    \newpage

    \silentsubsection{Вопрос 2}

    \textbf{Условие:}

    Доказать, что выполнение одновременно аксиомы Архимеда и леммы о стягивающихся отрезках
    равносильно выполнению аксиомы полноты.

    \textit{(1 балл)}

    \bigskip \bigskip

    \begin{itemize}
        \item
            \( \Rightarrow \)

            Есть \( A \leq B \)

            Изначально возьмем \( a_1 \in A, b_1 \in B \).

            Если так совпало, что \( a_1 = b_1 \), то можно взять \( c = a_1 \).

            И правда \( \forall a \in A, b \in B : a \leq b_1 = c, c = a_1 \leq b \)

            Иначе будем строить последовательность стягивающихся отрезков.

            \( I_1 = [a_1, b_1] \).

            Заметим, что в \( I_1 \) есть как элементы \( A \), так и элементы \( b \).

            Поделим \( I_1 \) пополам.
            Если в какой-то половине лежат числа и из \( A \), и из \( B \),
            то сделаем следующим отрезком ее.
            В противном случае разделителем является середина \( c = \frac{a_1 + b_1}{2} \):

            Чтобы это доказать, рассмотрим \( a \in A \):

            \( b_1 \in B, A \leq B \Rightarrow a \leq b_1 \),
            но \( \left[ \frac{a_1 + b_1}{2}, b_1 \right] \cap A = \varnothing \),
            а значит \( a \leq c \).

            Аналогично \( \forall b \in B \ c \leq b \)

            Продолжим построение остальных отрезков аналогично:

            Есть отрезок \( I_n = [a_n, b_n] \).
            Важно, что в этом отрезке содержатся элементы и \( A \), и \( B \).
            Если в какой-то из половин
            \(
                \left[ a_n, \frac{a_n + b_n}{2} \right), \left( \frac{a_n + b_n}{2}, b_n \right]
            \)
            есть числа и из \( A \), и из \( B \), то возьмем ее следующим отрезком.
            Иначе разделителем можно взять \( \frac{a_n + b_n}{2} \):

            Действительно, возьмем какое-нибудь число \( b \in B \cap I_n \)
            (такое обязательно есть, так как \( I_n \cap A, I_n \cap B \neq \varnothing \)).
            \( A \leq B \Rightarrow \forall a \in A \ a \leq b \leq b_n \),
            но \( \left( \frac{a_n + b_n}{2}, b_n \right] \cap A = \varnothing \),
            поэтому \( \forall a \in A \ a \leq \frac{a_n + b_n}{2} \).

            Аналогично \( \forall b \in B \ \frac{a_n + b_n}{2} \leq b \).

            Либо когда-то найдем разделитель, либо получим последовательность вложенных отрезков.

            По Бернулли \( 2^n = (1 + 1)^n \geq 1 + n \)

            \( | I_n | = \frac{b_1 - a_1}{2^{n - 1}} \leq \frac{1}{n} \)

            По аксиоме Архимеда \( \forall \varepsilon > 0 \ \exists n : \frac{1}{n} < \varepsilon \).

            Значит длина отрезка стремится к нулю \( \Rightarrow \{ I_n \}_{n = 1}^{+ \infty} \) ---
            последовательность стягивающихся отрезков.

            А у таких отрезков всегда есть общая точка \( c \).

            Предположим, что это не разделитель:

            Не умаляя общности, \( \exists a' \in A : c < a' \).
            Тогда \( a_n < c < a' \leq b_n \) (выше показывалось, почему \( A \leq \{ b_n \} \)).
            Получаем, что \( | I_n | = b_n - a_n > a' - c \),
            то есть последовательность отрезков не является стягивающейся --- противоречие.
        \item
            \( \Leftarrow \)

            \begin{enumerate}
                \item
                    Покажем выполнение леммы о стягивающихся отрезках

                    Возьмем \( A = \{ a_1, a_2, \ldots, a_n, \ldots \}, B = \{ b_1, b_2, \ldots, b_n, \ldots \} \).

                    Покажем, что \( A \leq B \):

                    Для этого покажем, что \( \forall i, j \in \NN \ a_i \leq b_j \):
                    \[
                        a_i \leq a_{ \max \{ i, j \} } \leq b_{ \max\{ i, j \} } \leq b_j
                    \]

                    (неравенства выше верны в силу вложенности отрезков)

                    Таким образом для любой последовательности вложенных отрезков у этих отрезков существует общая точка.

                    Если длины стремятся к \( 0 \), то такая точка одна, иначе:
                    \begin{gather*}
                        \forall n \in \NN \ a_n \leq c_1 < c_2 \leq b_n
                        \\
                        \Downarrow
                        \\
                        \forall n \in \NN \ | I_n | > c_2 - c_1
                    \end{gather*}

                    Противоречие с тем, что последовательность стягивающаяся.
                \item
                    Покажем выполнение аксиомы Архимеда

                    Покажем, что  \( \forall \varepsilon > 0 \ \forall x \ \exists n : n \varepsilon > x \).

                    Зафиксируем \( \varepsilon > 0 \)

                    Пусть \( A = \{ x \in \RR : \exists n : x < n \varepsilon \}, B = \RR \setminus A \)

                    Если \( B = \varnothing \), то лемма Архимеда верна. Предположим обратное.

                    \( 0 < \varepsilon \Rightarrow 0 \in A \).

                    Значит \( A, B \neq \varnothing \).

                    Предположим, что \( A \leq B \) неверно.
                    Тогда \( \exists a \in A, b \in B : b < a \).
                    Но
                    \(
                        a \in A \Rightarrow \exists n : a < n \varepsilon
                        \Rightarrow
                        b < a < n \varepsilon \Rightarrow b \in A \Rightarrow b \notin B
                    \) --- противоречие.

                    Значит существует разделитель \( c \).

                    Покажем, что разделитель единственный.
                    Пусть есть \( c_1 < c_2 \).
                    Но тогда \( \forall a \in A, b \in B \) выполнено \( a < c_1 < c_2 < b \),
                    а значит \( \frac{c_1 + c_2}{2} \notin A, B \), но \( B = \RR \setminus A \) --- противоречие.

                    Заметим, что \( A \leq \{ c \} \leq B \), а значит все числа \( < c \) принадлежат \( A \)
                    (так как каждое число лежит либо в \( A \), либо в \( B \)).

                    Но тогда рассмотрим \( c - \varepsilon \):
                    \begin{gather*}
                        c - \varepsilon < c \Rightarrow c - \varepsilon \in A
                        \\
                        \Downarrow
                        \\
                        \exists n \in \NN : c - \varepsilon < n \varepsilon \Rightarrow c + \varepsilon < (n + 2) \varepsilon
                        \\
                        \Downarrow
                        \\
                        c + \varepsilon \in A
                    \end{gather*}

                    Но \( c < c + \varepsilon \) --- противоречие.

                    Значит изначально было \( B = \varnothing \Rightarrow A = \RR \),
                    а именно это мы и хотели доказать.
            \end{enumerate}
    \end{itemize}

    \newpage

    \silentsection{Листик 3}

    \begin{center}
        \textit{(Последовательности и число \( e \))}
    \end{center}

    \silentsubsection{Вопрос 1}

    \textbf{Условие:}

    Пусть \( \lim\limits_{n \to \infty} a_n = a \)

    \silentsubsubsection{а)}

    \textbf{Условие:}

    \( \lim\limits_{n \to \infty} \frac{a_1 + a_2 + \ldots + a_n}{n} = a \)

    \textit{(1.5 балла)}

    \bigskip \bigskip

    Рассмотрим \( \varepsilon > 0 \)

    \( \exists K : \forall n > K \ |a_n - a| < \frac{\varepsilon}{2} \).

    Пусть \( C = a_1 + a_2 + \ldots + a_K \)
    \begin{gather*}
        \forall n > K
        \\
        \frac{C + (n - K)(a - \frac{\varepsilon}{2})}{n}
            < \frac{a_1 + a_2 + \ldots + a_n}{n}
            < \frac{C + (n - K)(a + \frac{\varepsilon}{2})}{n}
    \end{gather*}

    Очевидно, что левая и правая части неравенства стремятся к
    \( a - \frac{\varepsilon}{2} \) и \( a + \frac{\varepsilon}{2} \), соответственно.
    \begin{gather*}
        \exists N_1 > K : \forall n > N_1 \ \frac{C + (n - K)(a - \frac{\varepsilon}{2})}{n} > a - \varepsilon
        \\
        \exists N_2 > K : \forall n > N_2 \ \frac{C + (n - K)(a + \frac{\varepsilon}{2})}{n} < a + \varepsilon
        \\
        \Downarrow
        \\
        N = \max \{ N_1, N_2 \}
        \\
        \forall n > N \ a - \varepsilon < \frac{a_1 + a_2 + \ldots + a_n}{n} < a + \varepsilon
    \end{gather*}

    А значит \( \lim\limits_{n \to \infty} \frac{a_1 + a_2 + \ldots + a_n}{n} = a \)

    Что и требовалось доказать.

    \silentsubsubsection{б)}

    \textbf{Условие:}

    Если \( \forall n \in \NN \ a_n > 0 \), то и \( \lim\limits_{n \to \infty} \sqrt[n]{a_1 a_2 \cdot \ldots \cdot a_n} = a \)

    \textit{(1 балл)}

    \bigskip \bigskip

    \begin{enumerate}
        \item
            \( a > 0 \)

            Рассмотрим \( 0 < \varepsilon < a \)

            \( \exists K : \forall n > K \ 0 < a - \frac{\varepsilon}{2} < a_n < a + \frac{\varepsilon}{2} \).

            Пусть \( C = a_1 a_2 \cdot \ldots \cdot a_K \)
            \begin{gather*}
                \forall n > K
                \\
                \sqrt[n]{C \left( a - \frac{\varepsilon}{2} \right)^{n - K}}
                    < \sqrt[n]{a_1 a_2 \cdot \ldots \cdot a_n}
                    < \sqrt[n]{C \left( a + \frac{\varepsilon}{2} \right)^{n - K}}
            \end{gather*}

            Очевидно, что левая и правая части неравенства стремятся к
            \( a - \frac{\varepsilon}{2} \) и \( a + \frac{\varepsilon}{2} \), соответственно.
            \begin{gather*}
                \exists N_1 > K : \forall n > N_1 \ \sqrt[n]{C \left( a - \frac{\varepsilon}{2} \right)^{n - K}} > a - \varepsilon
                \\
                \exists N_2 > K : \forall n > N_2 \ \sqrt[n]{C \left( a + \frac{\varepsilon}{2} \right)^{n - K}} < a + \varepsilon
                \\
                \Downarrow
                \\
                N = \max \{ N_1, N_2 \}
                \\
                \forall n > N \ a - \varepsilon < \sqrt[n]{a_1 a_2 \cdot \ldots \cdot a_n} < a + \varepsilon
            \end{gather*}

            А значит \( \lim\limits_{n \to \infty} \sqrt[n]{a_1 a_2 \cdot \ldots \cdot a_n} = a \)
        \item
            \( a = 0 \)

            Рассмотрим \( 0 < \varepsilon \)

            \( \exists K : \forall n > K \ 0 < a_n < \frac{\varepsilon}{2} \).

            Пусть \( C = a_1 a_2 \cdot \ldots \cdot a_K \)
            \begin{gather*}
                \forall n > K
                \\
                0 < \sqrt[n]{a_1 a_2 \cdot \ldots \cdot a_n}
                    < \sqrt[n]{C \left( \frac{\varepsilon}{2} \right)^{n - K}}
            \end{gather*}

            Очевидно, что правая часть неравенства стремится к \( \frac{\varepsilon}{2} \)
            \begin{gather*}
                \exists N > K : \forall n > N \ \sqrt[n]{C \left( \frac{\varepsilon}{2} \right)^{n - K}} < \varepsilon
                \\
                \Downarrow
                \\
                \forall n > N \ 0 - \varepsilon < 0 < \sqrt[n]{a_1 a_2 \cdot \ldots \cdot a_n} < 0 + \varepsilon
            \end{gather*}

            А значит \( \lim\limits_{n \to \infty} \sqrt[n]{a_1 a_2 \cdot \ldots \cdot a_n} = 0 \)
    \end{enumerate}

    Что и требовалось доказать.

    \newpage

    \silentsubsubsection{в)}

    Если \( \forall n \in \NN \ a_n > 0 \) и \( \lim\limits_{n \to \infty} \frac{a_{n + 1}}{a_n} \),
    то \( \lim\limits_{n \to \infty} \sqrt[n]{a_n} = a \)

    \textit{(1 балл)}

    \bigskip \bigskip

    \begin{gather*}
        b_n = \frac{a_{n + 1}}{a_n}
        \\
        b_n > 0
        \\
        \lim\limits_{n \to \infty} b_n = a
        \\
        \lim\limits_{n \to \infty} \sqrt[n]{a_n} = \lim\limits_{n \to \infty} \sqrt[n] {b_1 b_2 \cdot \ldots \cdot b_{n - 1}}
    \end{gather*}

    Сводим к предыдущему пункту (внутри корня \( n - 1 \) множитель, а не \( n \), но не критично).

    \newpage

    \silentsubsection{Вопрос 2}

    \textbf{Условие:}

    \textbf{(Теорема Штольца).}
    Пусть \( \forall n \in \NN \ y_{n + 1} > y_n > 0 \) и \( \lim\limits_{n \to \infty} y_n = +\infty \),
    а также \( \lim\limits_{n \to \infty} \frac{x_{n + 1} - x_n}{y_{n + 1} - y_n} = l \).
    Доказать, что \( \lim\limits_{n \to \infty} \frac{x_n}{y_n} = l \).

    \textit{(1.5 балла, если выводится из теоремы Теплица и 2 балла, если доказывается независимо)}

    (Артем Александрович говорил, что то, как доказывать Штольца, будет решать проверяющий, а не Вы)

    \bigskip \bigskip

    \begin{itemize}
        \item
            Из Теплица:

            \begin{gather*}
                a_n = \frac{x_{n + 1} - x_n}{y_{n + 1} - y_n}
                \\
                c_{i j} = \frac{y_{j + 1} - y_j}{y_{i + 1}}
            \end{gather*}

            Убедимся, что \( c_{i j} \) удовлетворяют условиям теоремы Теплица:
            \begin{gather*}
                \lim\limits_{n \to \infty} c_{n k} = \frac{y_{k + 1} - y_k}{y_{n + 1}} = 0
                \\
                \lim\limits_{n \to \infty} \sum\limits_{k = 1}^{n} c_{n k}
                    = \lim\limits_{n \to \infty} \frac{y_{n + 1} - y_1}{y_{n + 1}}
                    = 1
                \\
                \sum\limits_{k = 1}^{n} |c_{n k}| = \frac{y_{n + 1} - y_1}{y_{n + 1}} < 1
            \end{gather*}

            \begin{gather*}
                b_n = \sum\limits_{k = 1}^{n} c_{nk} a_k
                    = \frac{x_2 - x_1}{y_2 - y_1} \cdot \frac{y_2 - y_1}{y_{n + 1}}
                    + \ldots
                    + \frac{x_{n + 1} - x_n}{y_{n + 1} - y_n} \cdot \frac{y_{n + 1} - y_n}{y_{n + 1}}
                    =
                \\
                = \frac{x_{n + 1} - x_1}{y_{n + 1}}
                \\
                \begin{rcases}
                    \lim\limits_{n \to \infty} b_n = l
                    \\
                    \lim\limits_{n \to \infty} \frac{x_1}{y_{n + 1}} = 0
                \end{rcases}
                \Rightarrow
                \lim\limits_{n \to \infty} \frac{x_{n + 1}}{y_{n + 1}} = l
            \end{gather*}

            Что и требовалось доказать.

            \newpage
        \item
            Независимо:

            Рассмотрим \( a_n = \frac{x_{n + 1} - x_n}{y_{n + 1} - y_n} \).
            После домножения на знаменатель:
            \begin{gather*}
                x_2 - x_1 = a_1 (y_2 - y_1)
                \\
                \vdots
                \\
                x_{n + 1} - x_n = a_n (y_{n + 1} - y_n)
            \end{gather*}

            Знаем, что \( a_n = l + \alpha_n \), где \( \alpha_n \) --- бесконечно малая последовательность.
            Сложим все равенства:
            \begin{gather*}
                x_{n + 1} - x_1 = l (y_{n + 1} - y_1) + \alpha_1 (y_2 - y_1) + \ldots + \alpha_n (y_{n + 1} - y_n)
                \\
                x_{n + 1} = l y_{n + 1} - l y_1 + \alpha_1 (y_2 - y_1) + \ldots + \alpha_n (y_{n + 1} - y_n) + x_1
                \\
                \frac{x_{n + 1}}{y_{n + 1}}
                    = l + \frac{\alpha_1 (y_2 - y_1) + \ldots + \alpha_n (y_{n + 1} - y_n) + x_1 - l y_1}{y_{n + 1}}
            \end{gather*}

            Давайте покажем, что правая часть стремится к \( l \).
            Для этого докажем, что дробь стремится к нулю.

            Очевидно, что \( \frac{x_1 - l y_1}{y_{n + 1}} \) стремится к \( 0 \)
            (в числителе константа, \( y_{n + 1} \to +\infty \)).
            Остается разобраться с альфами.

            Рассмотрим \( \varepsilon > 0 \)

            \( \exists K : \forall n > K \ 0 < \alpha_n < \frac{\varepsilon}{2} \).

            Пусть \( C = \alpha_1 (y_2 - y_1) + \ldots + \alpha_K (y_{K + 1} - y_K) \)
            \begin{gather*}
                b_n = \frac{\alpha_1 (y_2 - y_1) + \ldots + \alpha_n (y_{n + 1} - y_n)}{y_{n + 1}}
                \\
                \forall n > K
                \\
                b_n < \frac{C + \frac{\varepsilon}{2} ((y_{K + 2} - y_{K + 1}) + \ldots + (y_{n + 1} - y_n))}{y_{n + 1}}
                \\
                b_n < \frac{C + \frac{\varepsilon}{2} (y_{n + 1} - y_{K + 1})}{y_{n + 1}}
            \end{gather*}

            Очевидно, что правая часть неравенства стремится к \( \frac{\varepsilon}{2} \)
            \begin{gather*}
                \exists N > K : \forall n > N \ \frac{C + \frac{\varepsilon}{2} (y_{n + 1} - y_{K + 1})}{y_{n + 1}} < \varepsilon
                \\
                \Downarrow
                \\
                \forall n > N \ 0 - \varepsilon < 0 \leq b_n < 0 + \varepsilon
            \end{gather*}

            А значит \( \lim\limits_{n \to \infty} b_n = 0 \).

            Что и требовалось доказать.
    \end{itemize}

    \newpage

    \silentsubsection{Вопрос 3}

    \textbf{Условие:}

    Доказать, что

    \bigskip \bigskip

    P.S.
    \begin{itemize}
        \item
            Нормально сформулированные доказательства а) и б) взяты из первого тома Фихтенгольца
        \item
            в), г) на эту тему есть \href{https://www.youtube.com/watch?v=uXaELjpY9ps}{видео} на youtube-канале HMath
            (прикольный канал, всем советую)
    \end{itemize}

    \silentsubsubsection{а)}

    \textbf{Условие:}

    \( \lim\limits_{n \to \infty} ( 1 + \frac{1}{1!} + \frac{1}{2!} + \ldots + \frac{1}{n!} ) = e \)

    \textit{(1.5 балла)}

    \bigskip \bigskip

    Пусть \( a_n = (1 + \frac{1}{n})^n  \) и \( b_n = 1 + \frac{1}{1!} + \frac{1}{2!} + \ldots + \frac{1}{n!} \).

    \begin{gather*}
        a_n = \left( 1 + \frac{1}{n} \right)^n
        \\
        = 1 + \frac{1}{1!} \left( 1 - \frac{0}{n} \right)
            + \frac{1}{2!} \left( 1 - \frac{0}{n} \right) \left( 1 - \frac{1}{n} \right)
            + \ldots
            + \frac{1}{n!} \left (1 - \frac{0}{n} \right) \ldots \left( 1 - \frac{n - 1}{n} \right)
        \\
        \leq 1 + \frac{1}{1!} + \frac{1}{2!} + \ldots + \frac{1}{n!} = b_n
    \end{gather*}

    Теперь давайте покажем, что \( \forall n \in \NN \ b_n \leq e \).

    Зафиксируем какое-то \( k \in \NN \) и покажем, что \( b_k \leq e \)
    \begin{gather*}
        \forall n > k
        \\
        \left( 1 + \frac{1}{n} \right)^n
        > 1 + \frac{1}{1!} \left( 1 - \frac{0}{n} \right)
            + \ldots
            + \frac{1}{k!} \left (1 - \frac{0}{n} \right) \ldots \left( 1 - \frac{k - 1}{n} \right)
    \end{gather*}

    После предельного перехода в неравенстве:
    \[
        e \geq 1 + \frac{1}{1!} + \ldots + \frac{1}{k!} = b_k
    \]

    Получили, что \( \forall n \in \NN \ a_n \leq b_n \leq e \).

    Поскольку \( \lim\limits_{n \to \infty} a_n = e \), по теореме о зажатом пределе \( \lim\limits_{n \to \infty} b_n = e \).

    Что и требовалось доказать.

    \newpage

    \silentsubsubsection{б)}

    \textbf{Условие:}

    \( e =  1 + \frac{1}{1!} + \frac{1}{2!} + \ldots + \frac{1}{n!} + \frac{\alpha_n}{n \cdot n!} \),
    причем \( 0 < \alpha_n < 1 \).

    \textit{(1.5 балла)}

    \bigskip \bigskip

    \begin{gather*}
        y_{n + p} - y_n = \frac{1}{(n + 1)!} + \frac{1}{(n + 2)!} + \ldots + \frac{1}{m!}
        \\
        y_{n + p} - y_n = \frac{1}{(n + 1)!} \left( 1 + \frac{1}{n + 1} + \ldots + \frac{1}{(n + 1)(n + 2) \ldots (n + p - 1)} \right)
        \\
        y_{n + p} - y_n \leq \frac{1}{(n + 1)!} \left( 1 + \frac{1}{n + 1} + \ldots + \frac{1}{(n + 1)^{(p - 1)}} \right)
    \end{gather*}

    После предельного перехода и применения суммы убывающей геометрической прогрессии:
    \begin{gather*}
        e - y_n \leq \frac{1}{(n + 1)!} \cdot \frac{n + 2}{n + 1}
        \\
        0 \leq e - y_{n + 1} < e - y_n < \frac{1}{(n + 1)!} \cdot \frac{n + 2}{n + 1}
            < \frac{1}{(n + 1)!} \cdot \frac{n + 1}{n} = \frac{1}{n \cdot n!}
        \\
        0 < \frac{\alpha_n}{n \cdot n!} < \frac{1}{n \cdot n!}
    \end{gather*}

    А значит действительно верно, что \( 0 < \alpha_n < 1 \)

    \newpage

    \silentsubsubsection{в)}

    \textbf{Условие:}

    \( e \) --- иррациональное

    \textit{(1 балл)}

    \bigskip \bigskip

    Пусть это не так и \( e = \frac{m}{n} \)

    Уже знаем, что
    \( \frac{m}{n} = e = 1 + \frac{1}{1!} + \frac{1}{2!} + \ldots + \frac{1}{n!} + \frac{\alpha_n}{n \cdot n!} \)
    и \( 0 < \alpha_n < 1 \)

    Но после умножения на \( n! \) обеих сторон равенства слева будет целое число,
    а справа --- нет, так как \( \frac{\alpha_n}{n} \) --- нецелое (так как \( 0 < \alpha_n < 1 \)).

    Получили противоречие \( \Rightarrow e \) --- иррациональное.

    \newpage

    \silentsubsubsection{г)}

    \textbf{Условие:}

    \( \lim\limits_{n \to \infty} \sin (2 \pi e n!) = 0 \)

    \textit{(1 балл)}

    \bigskip \bigskip

    \begin{gather*}
        \sin (2 \pi e n!)
        =
        \\
        =
        \sin \left( 2 \pi n! \left( 1 + \frac{1}{1!} + \frac{1}{2!} + \ldots + \frac{1}{n!} + \frac{\alpha_n}{n \cdot n!} \right) \right)
        =
        \\
        =
        \sin \left(
            2 \pi \left( n! + \frac{n!}{1!} + \frac{n!}{2!} + \ldots + \frac{n!}{n!} \right) + \frac{2 \pi \alpha_n}{n}
            \right)
        =
        \\
        =
        \sin \left( \frac{2 \pi \alpha_n}{n} \right)
    \end{gather*}


    Поскольку \( 0 < \alpha_n < 1 \), очевидно, что \( \lim\limits_{n \to \infty} \frac{2 \pi \alpha_n}{n} = 0 \).

    По непрерывности синуса (если ее нельзя, то лапками разбирайтесь)
    \[
        \lim\limits_{n \to \infty} \sin (2 \pi e n!)
        =
        \lim\limits_{n \to \infty} \sin \left( \frac{2 \pi \alpha_n}{n} \right)
        =
        0
    \]

    \newpage

    \silentsection{Листик 4}

    \begin{center}
        \textit{(Частичные пределы и критерий Коши)}
    \end{center}

    \silentsubsection{Вопрос 1}

    \textbf{Условие:}

    Доказать, что из последовательности функций \( f_n : A \to [0, 1] \) можно выбрать подпоследовательность,
    сходящуюся в каждой точке \( A \), если:

    \silentsubsubsection{а)}

    \textbf{Условие:}

    \( A = \{ 1, 2, \ldots, n \} \)

    \textit{(1 балл)}

    \bigskip \bigskip

    Поступим так же, как и со сходящейся подпоследовательностью.

    Но теперь будем опираться на \( n \) отрезков --- для каждого элемента \( A \).

    Пусть у нас есть набор отрезков \( ( J_1, J_2, \ldots, J_n ) \).
    Тогда функция \( f \) удовлетворяет набору, если  \( \forall i \in \{ 1, 2, \ldots, n \} \ f(i) \in J_i \)

    Изначально для каждого у нас набор вида \( J_1 = J_2 = \ldots = J_n = [0, 1] \).

    Теперь давайте поделим каждый отрезок пополам.
    Теперь сделаем \( 2^n \) наборов из половин (то есть возьмем от каждого отрезка только половину).
    Причем \( i \)-й элемент набора (\( J_i \)) -- какая-то половина отрезка для \( f(i) \).
    Давайте для каждого набора возьмем все такие функции, которые ему удовлетворяют.

    Заметим, что поскольку наборов конечно, а набор, из которого мы строили новые, удовлетворяет счетное число функций,
    какому-то набору половин удовлетворяет счетное число функций.

    Давайте возьмем его новым набором и продолжим шаг.

    Делая так, мы построим последовательности стягивающихся отрезков сразу для \( f(1), f(2), \ldots, f(n) \).

    (понятно, что нужно еще выбрать кандидатов из функций, которые удовлетворяют каждому набору, но это несложно,
    так как на каждом шаге у нас будет ``заблокировано'' лишь конечное число функций)

    \newpage

    \silentsubsubsection{б)}

    \textbf{Условие:}

    \( A = \NN \)

    \textit{(1.5 балла)}

    \bigskip \bigskip

    Нам не нужна равномерная сходимость, поэтому давайте немного адаптируем решение из пункта а).

    Будем называть удовлетворяющим набору размера \( n \) все функции такие,
    что \( \forall i \in \{ 1, 2, \ldots, n \} \ f(i) \in J_i \)
    (то есть теперь нас интересует только префикс значений области определения, потому что набор конечный).

    Пусть у нас сейчас набор \( ( J_1, J_2, \ldots, J_n ) \).

    Немного изменим процесс выбора следующего набора.
    Помимо половин, давайте просто ко всем наборам добавим \( J_{n + 1} = [0, 1] \) --- стартовый отрезок для \( f(n + 1) \).

    То есть после каждого шага наш набор будет увеличиваться на \( 1 \) отрезок.

    Продолжая строить так отрезки, мы сможем построить последовательность стягивающихся отрезков для любого \( f(n) \).

    \newpage

    \silentsubsection{Вопрос 2}

    \textbf{Условие:}

    Найти все частичные пределы, точные верхние и точные нижние грани последовательности \( a_n \)

    \silentsubsubsection{а)}

    \textbf{Условие:}

    \( a_n = \sin n \)

    \textit{(1.5 балла)}

    \bigskip \bigskip

    Кажется, это называется теоремой Кронекера.

    Покажем, что
    \begin{gather*}
        \forall \varepsilon > 0, \ a \in [0, 2 \pi), \ N \in \NN
        \\
        \exists n, k \in \NN : n > N, \ |2 \pi k + a - n| < \varepsilon
    \end{gather*}

    Зафиксируем \( \varepsilon > 0 \).

    Теперь разобьем окружность длиной \( 2 \pi \) на дуги размера \( \varepsilon \) (последняя дуга, возможно будет меньше).

    Начнем отмечать, куда попадают \( n > N \) на этой окружности.
    По принципу Дирихле в какую-то дугу точно попадут две точки.
    Если они совпали, то \( m - n = 2 \pi k \), где \( n, m, k \in \NN \) --- противоречие, ведь \( \pi \) иррациональное.

    В противном случае мы знаем, что за \( m - n \) ``прыжков'' точки мы умеем сдвигаться по окружности \( \leq \varepsilon \)
    в какую-то сторону.
    Давайте теперь совершать толькор такие прыжки --- поскольку их длина не превышает длины дуги,
    мы обязательно попадем во все дуги, на которые мы разбили окружность.
    А значит можем подобраться к любому числу на этой окружности с любой точностью
    (причем отбросив первые члены последовательности).

    Отсюда понятно, что точные нижняя и верхняя грани \( a_n \) --- \( -1 \) и \( 1 \), соответственно,
    а все частичные пределы --- \( [-1, 1] \)

    \textbf{Ответ:}

    \( \inf a_n = -1 \)

    \( \sup a_n = 1 \)

    Все частичные пределы --- \( [-1, 1] \)

    \newpage

    \silentsubsubsection{б)}

    \textbf{Условие:}

    \( a_n = \max \{ |\sin n|, |\sin (n + 1)| \} \)

    \textit{(1.5 балла)}

    \bigskip \bigskip

    Давайте поймем, что \( a_n \) может быть бесконечно близко к любому числу на отрезке \( [\sin 0.5, 1] \)

    И правда, пусть надо приблизить
    \[
        \sin x, x \in \left[ 0.5, \frac{\pi}{2} \right]
    \]

    Мы точно можем найти \( \forall N \in \NN \ n > N \) такое,
    что оно на окружности попадет в \( (x, x + \varepsilon) \) или \( (x - \varepsilon, x) \)
    (главное, чтобы этот интервал был внутри \( \left[ 0.5, \frac{\pi}{2} \right] \))
    для любого достаточно маленького \( \varepsilon > 0 \).

    Заметим, что в таком случае,
    \begin{gather*}
        0 \leq |n - 1| \leq |n| < \frac{\pi}{2}
        \\
        \Downarrow
        \\
        \max \{ |\sin (n - 1)|, |\sin n| \} = \sin n
    \end{gather*}

    В силу непрерывности синуса получим, что  \( \sin n \) может достаточно точно приблизить \( \sin x \).

    С другой стороны понятно, что если \( | \sin n | < \sin 0.5 \), то \( | \sin (n + 1) | \geq \sin 0.5 \), и наоборот.
    А значит все значения \( \{ a_n \} \) хотя бы \( \sin 0.5 \).

    Тогда понятно, что нижняя и верхняя грани --- \( \sin 0.5, \) и \( 1 \),
    а \( [ \sin 0.5, 1 ] \) --- все частичные пределы.

    \textbf{Ответ:}

    \( \inf a_n = \sin 0.5 \)

    \(\sup a_n = 1 \)

    Все частичные пределы --- \( [\sin 0.5, 1] \)

    \newpage

    \silentsubsubsection{в)}

    \textbf{Условие:}

    Пусть \( s_n = a_1 + a_2 + \ldots + a_n \to +\infty \), при любом \( k \) выполнено \( a_k > 0 \),
    \( \lim\limits_{k \to \infty} a_k = 0 \).

    Найти множество частичных пределов дробных частей элементов \( s_n \).

    \textit{(1.5 балла)}

    \bigskip \bigskip

    Делаем то же самое, что и в предыдущих пунктах, только теперь окружность длины \( 1 \).

    Разобьем ее на дуги длины \( \varepsilon \).
    Поскольку \( a_n > 0 \) и \( \lim\limits_{n \to \infty} a_n = 0 \), с какого-то момента \( 0 < a_n < \varepsilon \).
    А значит прыжки по окружности становятся меньше длины дуги.
    \( \lim\limits_{n \to \infty} s_n = +\infty \), поэтому нам хватит суммарной длины прыжков,
    чтобы попасть в каждую дугу.

    Таким образом, мы умеем попадать в любую дугу заданного размера.
    Значит \( s_n \) может быть бесконечно близко к любой точке окружности.
    Отсюда искомые частичные пределы --- \( [0, 1) \).

    \textbf{Ответ:}

    \( [0, 1) \)

    \newpage

    \silentsubsection{Вопрос 3}

    \textbf{Условие:}

    Имеет ли предел последовательность \( a_n = \sum\limits_{k = 1}^{n} \frac{\cos \sqrt{k}}{\sqrt{k}} \) ?

    \textit{(1.5 балла)}

    \bigskip \bigskip

    Покажем невыполнение критерия Коши.

    Постараемся найти много \( \cos \sqrt{k} \geq \frac{1}{2} \), идущих подряд.

    \( 2 \pi l - \frac{\pi}{3} \leq \sqrt {k} \leq 2 \pi l + \frac{\pi}{3}, \ l \in \NN \) --- в таком случае
    косинус будет именно таким, каким нам нужно. Какие тогда ограничения на \( k \)?
    \begin{gather*}
        2 \pi l - \frac{\pi}{3} \leq \sqrt {k} \leq 2 \pi l + \frac{\pi}{3}, \ l \in \ZZ
        \\
        \left( 2 \pi l - \frac{\pi}{3} \right)^2 \leq k \leq \left( 2 \pi l + \frac{\pi}{3} \right)^2
        \\
        4 \pi^2 l^2 - \frac{4 \pi^2 l}{3} + \frac{\pi^2}{9} \leq k \leq 4 \pi^2 l^2 + \frac{4 \pi^2 l}{3} + \frac{\pi^2}{9}
    \end{gather*}

    Получили отрезок прямой длиной \( \frac{8 \pi^2 l}{3} \).

    Очевидно, что в нем хотя бы \( \frac{8 \pi^2 l}{3} - 1 \) натуральных чисел.

    Пусть это числа от \( L \) до \( R \).

    Но тогда
    \begin{gather*}
        \sum\limits_{k = L}^{R} \frac{\cos \sqrt{k}}{\sqrt{k}} \geq \sum\limits_{k = L}^{R} \frac{1}{2 \sqrt k}
            \geq \frac{R - L + 1}{2 \sqrt R}
        \\
        \sum\limits_{k = L}^{R} \frac{\cos \sqrt{k}}{\sqrt{k}}
            \geq \frac{\frac{8 \pi^2 l}{3} - 1}{2 \pi l + \frac{\pi}{3}}
            \geq \frac{\pi^2 l}{4 \pi l}
            = \frac{\pi}{4}
    \end{gather*}

    Поскольку \( \forall N \in \NN \) мы можем найти пару \( (L, R) \) такую, что \( L > N \),
    получили противоречие критерию Коши.

    \newpage

    \silentsection{Листик 5}

    Очень подозрительно, но тут нет дополнительных вопросов к коллоквиуму...

    \newpage

    \silentsection{Листик 6}

    много буков не осилил

    \newpage
\end{document}
